\section{M8 Measure Localization Route}
\label{sec:m8_measure_localization_route}
\begingroup
\sloppy
\emergencystretch=3em

This fragment records the M8 route from the current represented-integral
global theorem to a measure-facing Stokes statement.  The goal is to make the
remaining analytic handoff explicit: indicators localize finite sums, the bulk
integrand is identified almost everywhere with the localized exterior
derivative integrand, the local scalar integrands are integrable, the boundary
measure is reconstructed from selected boundary pieces, and the M8 statement
packages those facts without redesigning the M7 constructor layer.

\begin{remark}[M8 graph conventions]
  \label{anc:m8_graph_conventions}
  \uses{anc:m7_parallel_worker_queue}
  Figure~\ref{fig:stokes_m8_measure_graph} separates three statuses that were
  previously easy to conflate, and
  Figure~\ref{fig:stokes_m8_proof_wave_graph} zooms in on the next true-proof
  wave.  Green nodes are existing Lean anchors or constructor packages.  Orange
  nodes are not vague future work: each one names a specific mathematical
  instance that still has to be produced for the existing fields.  Blue nodes
  are theorem targets or final wrappers that should become short once the
  orange nodes are discharged.
\end{remark}

\begin{theorem}[M8 indicator localization theorem]
  \label{anc:m8_indicator_theorem}
  \lean{Stokes.indicator_eq_self_of_support_subset,
    Stokes.ae_eq_indicator_of_support_subset,
    Stokes.finset_sum_ae_eq_indicator_sum_of_support_subset,
    Stokes.Icc_indicator_finset_sum_ae_eq_sum_of_support_subset,
    Stokes.integral_eq_finset_sum_setIntegral_of_ae_eq_sum_indicator,
    Stokes.integral_eq_finset_sum_setIntegral_of_support_subset}
  \leanok
  \uses{anc:m6_4_partition_localization_support_control,
    anc:m6_5_global_integral_extderiv_reconstruction}
  \emph{Lean modules:}
  \texttt{Stokes.Global.IndicatorSupportLocalization} and
  \texttt{Stokes.Global.MeasureIntegralLocalization}.

  \emph{Role in M8:}
  this is the pure measure-theory bridge from support control to finite
  additivity.  Termwise support containment gives pointwise and a.e. equality
  between an unlocalized finite sum and the corresponding sum of indicator
  terms.  The Bochner-integral theorem then turns an a.e. equality
  \[
    F = \sum_{i \in A} \mathbf{1}_{K_i} f_i
  \]
  into
  \[
    \int F\,d\mu =
      \sum_{i \in A} \int_{K_i} f_i\,d\mu,
  \]
  provided the localized terms are integrable.

  \emph{M8 handoff:}
  later workers should not reprove finite-sum integral algebra inside the bulk
  or boundary packages.  They should instead prove the support containment,
  measurability, integrability, and a.e. equality hypotheses needed to call
  these indicator theorems.
\end{theorem}

\begin{theorem}[M8 bulk a.e. equality]
  \label{anc:m8_bulk_ae_equality}
  \lean{Stokes.bulkIntegrand,
    Stokes.bulkIntegrand_apply_eq_of_extDeriv_eq,
    Stokes.bulkIntegrand_ae_eq_of_forall_extDeriv_eq,
    Stokes.BulkIntegrandAEData,
    Stokes.BulkIntegrandAEData.ofExtDerivOnSupportData,
    Stokes.BulkIntegrandAEData.ofExtDerivPartitionReconstructionData,
    Stokes.BulkIntegrandAEData.ofLocalizedFormEventuallyEqData}
  \leanok
  \uses{anc:m8_indicator_theorem,
    anc:m6_4_partition_localization_support_control,
    anc:m6_5_global_integral_extderiv_reconstruction}
  \emph{Lean module:} \texttt{Stokes.Global.BulkIntegrandAE}.

  \emph{Role in M8:}
  this node identifies the scalar top-degree bulk integrand for
  \(\dd(\sum_i \rho_i\omega)\) with the scalar top-degree integrand for
  \(\dd\omega\) for the chartwise measures used by the bulk integral.  The
  existing package is deliberately measure-parametric: a worker supplies the
  chartwise measure and proves that its a.e. support lies in the chart-support
  set where the exterior-derivative reconstruction is valid.

  \emph{Explicit field mathematics:}
  the next bulk-localization theorem still has to combine this a.e. equality
  with the indicator theorem above.  Its concrete output should be a genuine
  measure integral equality for the represented bulk side, not another global
  Stokes constructor.  Proposed Lean names:
  \plannedlean{Stokes.bulkMeasureIntegral_eq_indicatorLocalBulkSum}
  \plannedlean{Stokes.bulkMeasureLocalizationData_of_bulkIntegrandAE}
\end{theorem}

\begin{theorem}[M8 integrability package]
  \label{anc:m8_integrability_package}
  \lean{Stokes.integrableOn_of_continuousOn_isCompact,
    Stokes.integrable_of_continuousOn_isCompact_support_subset,
    Stokes.integrableOn_Icc_of_contDiffOn,
    Stokes.LocalizedSmoothnessData.integrableOn_Icc,
    Stokes.BoundaryIntegrabilityData,
    Stokes.BoundaryIntegrabilityData.of_extendedBox_chartwiseSmooth_targetBox,
    Stokes.boundaryChartTransitionPullbackIntegrand_integrableOn_of_extendedBox,
    Stokes.boundaryChartInChartIntegrand_integrableOn_of_chartwiseSmooth_selectedBox}
  \leanok
  \uses{anc:m6_4_active_compact_support_to_boxes,
    anc:m6_4_partition_localization_support_control,
    anc:m6_4_boundary_pieces_chart_change_fields,
    anc:m8_indicator_theorem}
  \emph{Lean modules:}
  \texttt{Stokes.Global.CompactSupportIntegrability} and
  \texttt{Stokes.Global.BoundaryIntegrabilityCompactSupport}.

  \emph{Role in M8:}
  the indicator and finite-sum integral theorems require integrability of every
  selected local term.  The compact-support wrappers prove the bulk side from
  continuity or \texttt{ContDiffOn} control on compact coordinate boxes.  The
  boundary wrappers prove the lower-zero-face integrability needed for source
  transition-pullback terms, Jacobian-weighted chart-change terms, and target
  in-chart boundary terms.

  \emph{Explicit field mathematics:}
  M8 workers should instantiate these wrappers with the exact selected boxes
  and lower-zero face boxes chosen by M6.4--M7.  The important alignment fields
  are the selected active set, the selected partition coefficient, the source
  and target box records, and the maps from those boxes to the support sets
  used by the actual measures.
\end{theorem}

\begin{theorem}[M8 boundary measure localization]
  \label{anc:m8_boundary_measure_localization}
  \lean{Stokes.boundaryMeasurePieceIndicator,
    Stokes.boundaryMeasurePieceSum,
    Stokes.boundaryMeasureIndicatorSum,
    Stokes.boundaryMeasureIntegral_eq_selectedBoundaryPieceSum_of_ae_eq,
    Stokes.boundaryMeasureIntegral_eq_selectedBoundaryPieceSum_of_ae_indicator_eq,
    Stokes.BoundaryMeasureLocalizationData,
    Stokes.BoundaryMeasureLocalizationData.boundaryMeasureIntegral_eq_selectedBoundaryPieceSum,
    Stokes.BoundaryMeasureLocalizationData.toBoundaryIntegralReconstructionData,
    Stokes.BoundaryMeasureLocalizationData.manifoldBoundaryIntegral_eq_selectedBoundaryPieceSum}
  \leanok
  \uses{anc:m8_indicator_theorem,
    anc:m8_integrability_package,
    anc:m6_4_boundary_pieces_chart_change_fields}
  \emph{Lean modules:}
  \texttt{Stokes.Global.BoundaryMeasureLocalization} and
  \texttt{Stokes.Global.BoundaryCOVMeasureConstructor}.

  \emph{Role in M8:}
  this node fieldizes the genuine boundary measure integral.  It takes an
  ambient boundary measure, an ambient boundary integrand, selected boundary
  piece sets, selected piece integrands, and the comparison from each piece
  integral to the recorded boundary partition term.  The result is the exact
  finite boundary partition sum consumed by the existing global reconstruction
  layer.

  \emph{Explicit field mathematics:}
  the boundary worker must prove the a.e. decomposition of the boundary
  integrand into selected indicator-localized pieces, termwise integrability,
  and the equality between every selected piece integral and the transported
  \texttt{boundaryPartitionTerm}.  Chart-change and COV work should feed this
  node through term equalities, not by changing the M8 statement shape.
\end{theorem}

\begin{theorem}[M8 measure-level Stokes statement]
  \label{anc:m8_measure_level_statement}
  \lean{Stokes.M8MeasureLocalizationData,
    Stokes.M8MeasureLocalizationData.toBulkIntegralReconstructionData,
    Stokes.M8MeasureLocalizationData.toBulkMeasureLocalizationFields,
    Stokes.M8MeasureLocalizationData.toBoundaryMeasureLocalizationFields,
    Stokes.M8GlobalStokesInput,
    Stokes.M8GlobalStokesInput.toNaturalGlobalStokesInput,
    Stokes.M8GlobalStokesInput.toNaturalMeasureStokesInput,
    Stokes.M8GlobalStokesInput.stokes,
    Stokes.m8GlobalStokes,
    Stokes.m8GlobalStokes_represented}
  \leanok
  \uses{anc:m7_selected_mixed_global_assembly,
    anc:m8_bulk_ae_equality,
    anc:m8_boundary_measure_localization,
    anc:m8_integrability_package}
  \emph{Lean module:} \texttt{Stokes.Global.M8Statement}.

  \emph{Role in M8:}
  this is the blueprint-facing statement node.  The input keeps the existing
  M7 mixed-constructor data visible and adds a single measure-localization
  package containing the genuine bulk and boundary measure integrals.  The
  proof projects the data to \texttt{NaturalGlobalStokesInput} and then to
  \texttt{NaturalMeasureStokesInput}; no new final constructor layer is needed.

  \emph{Remaining route:}
  the statement already exists as bookkeeping.  The mathematical work is to
  instantiate \texttt{M8MeasureLocalizationData}: prove bulk measure
  localization from the bulk a.e. equality and integrability package, prove
  boundary measure localization from the boundary measure node, and align both
  with the same selected active set, target-image family, artificial-face
  cancellation, and chart-change cancellation used in M7.  Proposed Lean name:
  \plannedlean{Stokes.m8MeasureLocalizationData_of_selectedPartition}
\end{theorem}

\begin{theorem}[M8 measure localization worker queue]
  \label{anc:m8_measure_localization_worker_queue}
  \uses{anc:m8_indicator_theorem,
    anc:m8_bulk_ae_equality,
    anc:m8_integrability_package,
    anc:m8_boundary_measure_localization,
    anc:m8_measure_level_statement}
  The measure-localization work can be split without touching the final
  constructor API:
  \begin{enumerate}
    \item prove the selected-box and selected-face support containments needed
      by the indicator theorem;
    \item prove the chartwise bulk a.e. equality for the measures used by the
      bulk integral;
    \item instantiate bulk and boundary integrability on every selected box or
      lower-zero face box;
    \item prove boundary a.e. decomposition and piece-integral comparisons;
    \item fill \texttt{M8MeasureLocalizationData} and call
      \texttt{m8GlobalStokes}.
  \end{enumerate}
  The queue deliberately leaves the existing M6.4--M7 records in charge of
  compact support, local pieces, artificial-face cancellation, and chart-change
  cancellation.
\end{theorem}

\begin{proposition}[Current twenty-agent M8 split]
  \label{anc:m8_current_twenty_agent_split}
  \uses{anc:m8_measure_localization_worker_queue,
    anc:m8_measure_level_statement}
  \begingroup\raggedright
  The current parallel wave should be read as a field-owner map, not as twenty
  independent final theorems.  The useful split is:
  \begin{enumerate}
    \item Bulk a.e. equality:
      \texttt{BulkIntegrandAE} under \texttt{Stokes.Global}.
    \item Indicator support localization:
      \texttt{IndicatorSupportLocalization} under \texttt{Stokes.Global}.
    \item Bulk localization constructor:
      \texttt{BulkIntegralLocalizationConstructor} under
      \texttt{Stokes.Global}.
    \item Boundary measure localization:
      \texttt{BoundaryMeasureLocalization} under \texttt{Stokes.Global}.
    \item Boundary COV measure constructor:
      \texttt{BoundaryCOVMeasureConstructor} under \texttt{Stokes.Global}.
    \item Boundary compact-support integrability:
      \texttt{BoundaryIntegrabilityCompactSupport} under
      \texttt{Stokes.Global}.
    \item Global integral definitions:
      \texttt{GlobalIntegralDefinitions} under \texttt{Stokes.Global}.
    \item Natural measure constructor:
      \texttt{NaturalMeasureConstructor} under \texttt{Stokes.Global}.
    \item M8 statement alignment:
      \texttt{M8Statement} under \texttt{Stokes.Global}.
    \item Artificial-face field reduction:
      \texttt{ArtificialFaceFieldReduction} under \texttt{Stokes.Global}.
    \item Pure boundary target-image reduction:
      \texttt{TargetImageFieldReduction} under
      \texttt{Stokes.BoundaryChart}.
    \item Target-image to global assembly adapter:
      \texttt{BoundaryTargetImageToAssembly} under \texttt{Stokes.Global}.
    \item Measure box API:
      \texttt{MeasureBoxAPI} under \texttt{Stokes.Global}.
    \item Bulk measure localization fields:
      \texttt{BulkMeasureLocalizationFields} under \texttt{Stokes.Global}.
    \item Measure-localization audit:
      \texttt{MeasureLocalizationAudit} under \texttt{Stokes.Global}.
    \item Focused build and import-graph audit.
    \item Remaining \texttt{FIELD} audit and blocker ledger.
    \item Blueprint and M8 dependency graph maintenance.
    \item Field-reduction report:
      \texttt{reports/stokes-m8-field-reduction.md}.
    \item Measure-route blueprint report and printable PDF.
  \end{enumerate}
  The parent integration rule is: keep pure boundary-chart modules out of
  \texttt{Stokes.HalfSpace}, keep global adapters in \texttt{Stokes.Global},
  and only promote modules to aggregators after their focused builds and the
  no-placeholder scan are clean.
  \par\endgroup
\end{proposition}

\begin{proposition}[Current twelve-agent M8 implementation wave]
  \label{anc:m8_current_twelve_agent_wave}
  \uses{anc:m8_measure_localization_worker_queue,
    anc:m8_p0_p1_p2_blocker_ledger}
  \begingroup\raggedright
  The current implementation wave is intentionally smaller than the previous
  twenty-agent restart.  It should focus on discharging true M8 instances
  rather than adding another layer of wrappers.  The twelve owners are:
  \begin{enumerate}
    \item bulk measure-localization fields from selected chart data;
    \item actual bulk a.e. integrand equality for the chartwise measure;
    \item measure-box comparison between measure-local and project-local terms;
    \item bulk localization constructor feeding \texttt{M8MeasureLocalizationData};
    \item ambient boundary measure and selected-piece indicator decomposition;
    \item boundary piece integrability and COV-to-partition term equalities;
    \item target boundary boxes, image equalities, and local inverse data;
    \item artificial-face support-zero or overlap-pairing geometry;
    \item boundary orientation bridge from oriented atlas/manifold data;
    \item compact-support selected interior and boundary box construction;
    \item final M8 assembly constructor calling \texttt{m8GlobalStokes};
    \item blueprint and report synchronization.
  \end{enumerate}
  The intended success condition is not twenty or twelve new theorem names; it
  is a smaller number of fields in \texttt{M8MeasureLocalizationData} whose
  values are produced from compact support, selected charts, measures, and
  orientation data.
  \par\endgroup
\end{proposition}

\begin{proposition}[Current eighteen-agent natural compact-support wave]
  \label{anc:m8_current_eighteen_agent_wave}
  \uses{anc:m8_current_twelve_agent_wave,
    anc:m8_to_manifold_dependency_ladder,
    anc:m8_measure_level_statement}
  \begingroup\raggedright
  Figure~\ref{fig:stokes_m8_measure_graph} now includes the current
  eighteen-agent wave as the explicit bridge from the conditional M8 theorem
  to a natural compact-support Stokes statement.  The owners are:
  \begin{enumerate}
    \item audit the shortest constructor route from M8 fields to natural
      compact-support inputs;
    \item define the compact-support-facing input record without changing the
      M8 theorem shape;
    \item construct finite selected chart boxes from compact support and chart
      domains;
    \item project selected partitions to the localized-interior fields used by
      M8;
    \item build the compact-support bulk measure package;
    \item adapt the bulk measure package to \texttt{M8MeasureLocalizationData};
    \item build the compact-support boundary measure package;
    \item adapt boundary measure reconstruction to the M8 boundary field;
    \item construct target-image data from local inverse and compact-image
      selections;
    \item adapt target-image data to the global M8 assembly fields;
    \item feed artificial-face support-zero or pairing data into the M8
      artificial-face fields;
    \item connect the boundary-chart sign convention to the induced oriented
      boundary orientation;
    \item align represented global integrals with the canonical compact-support
      integral interface;
    \item state the compact-support-facing M8 theorem as a short wrapper around
      \texttt{m8GlobalStokes};
    \item assemble a natural compact-support input constructor from the pieces
      above;
    \item integrate the new modules serially through aggregators after focused
      checks;
    \item maintain this blueprint graph and printable PDF without running Lean;
    \item perform the final parent-side serialized build and no-placeholder
      scan.
  \end{enumerate}
  The mathematical bottlenecks are still Agents 3, 12, and 13: finite chart-box
  selection from compact support, induced boundary orientation, and the
  canonical integral comparison.  The other owners mostly reduce field surface
  area by turning manual M8 fields into named constructors.
  \par\endgroup
\end{proposition}

\begin{proposition}[Proof-wave true-proof nodes]
  \label{anc:m8_proof_wave_true_nodes}
  \uses{anc:m8_current_eighteen_agent_wave,
    anc:m8_bulk_ae_equality,
    anc:m8_integrability_package,
    anc:m8_boundary_measure_localization,
    anc:m8_measure_level_statement}
  \begingroup\raggedright
  The next proof wave should not add another layer of M8 wrappers.  Its purpose
  is to turn the orange mathematical inputs in
  Figure~\ref{fig:stokes_m8_measure_graph} into theorem-produced data.
  Figure~\ref{fig:stokes_m8_proof_wave_graph} isolates this proof wave so that
  workers can take one node without editing the M8 assembly wrapper.  The
  following nodes are the intended isolated proof tasks.

  \begin{enumerate}
    \item \textbf{Bulk a.e. theorem.}  Prove the selected-partition version of
      \texttt{BulkIntegrandAEData}: for the chartwise bulk measures actually
      used by the global integral, the scalar integrand of
      \(\dd(\rho_i\omega)\) agrees a.e. with the localized exterior-derivative
      integrand on the support box.  This is the true proof behind
      \texttt{ExtDerivToBulkMeasure}; it should output the a.e. input consumed
      by \texttt{CompactSupportBulkMeasure}.

    \item \textbf{MeasureBoxAPI comparison.}  Prove that the measure-local bulk
      term for each selected chart is definitionally or theorem-equal to the
      project-local box term produced by the half-space/local Stokes package.
      This is the equality that lets the bulk finite sum in
      \texttt{BulkIntegralLocalizationConstructor} be the same finite sum that
      M8 receives from local Stokes.

    \item \textbf{Boundary measure theorem.}  Instantiate
      \texttt{BoundaryMeasureLocalizationData} from the actual boundary measure,
      selected boundary-piece sets, boundary integrands, integrability on lower
      zero faces, and the a.e. indicator decomposition of the boundary
      integrand.  The output should be
      \texttt{boundaryMeasureIntegral\_eq\_partitionSum}, not a new statement
      of Stokes.

    \item \textbf{Artificial support-zero route.}  For compactly supported
      localized forms whose support lies strictly inside the selected interior
      chart box away from artificial faces, prove that every artificial face
      contribution in the local box theorem is zero.  This feeds
      \texttt{ArtificialFaceSupportZeroToM8}.

    \item \textbf{Artificial adjacency route.}  In parallel with the
      support-zero route, construct the adjacent-face pairing for selected
      interior boxes: opposite orientation, the same geometric face, and equal
      unsigned face term.  This feeds \texttt{ArtificialFacePairingToM8}.  Only
      one of the support-zero or adjacency routes is needed for a first natural
      theorem, but both are legitimate proof targets.

    \item \textbf{Boundary orientation bridge.}  Connect the project
      \texttt{BoundaryChartOrientedAtlas} predicate to the intended mathlib
      oriented-manifold data and induced outward-normal-first boundary
      orientation.  This should justify the sign used by the lower-zero-face
      local Stokes theorem across boundary chart changes.
  \end{enumerate}

  These six nodes are the current true-proof frontier.  The surrounding M8
  input builders, compact-support statements, target-image adapters, and
  canonical-integral skeletons are useful only insofar as they expose these
  exact proof obligations cleanly.
  \par\endgroup
\end{proposition}

\begin{remark}[Focused-build policy for parallel M8 waves]
  \label{anc:m8_twelve_agent_focused_build_policy}
  \uses{anc:m8_current_twelve_agent_wave,
    anc:m8_current_eighteen_agent_wave,
    anc:m8_proof_wave_true_nodes}
  \begingroup\raggedright
  Parallel Lean workers should avoid running a full \texttt{lake build}.  Each
  worker should build only the narrowest target that covers its files, such as
  \texttt{lake build Stokes.Global.BulkIntegrandAE} or
  \texttt{lake build Stokes.BoundaryChart.TargetImageFieldReduction}.  The
  parent integration step should serialize the final full build only after
  workers are complete and aggregator imports are stable.  Documentation-only
  workers should not run Lean builds at all.  This is a build-engine constraint,
  not a mathematical shortcut: it avoids concurrent \texttt{.lake/build}
  writers and keeps failures local to the responsible owner.
  \par\endgroup
\end{remark}

\begin{remark}[M8 P0/P1/P2 blocker ledger]
  \label{anc:m8_p0_p1_p2_blocker_ledger}
  \uses{anc:m8_bulk_ae_equality,
    anc:m8_integrability_package,
    anc:m8_boundary_measure_localization,
    anc:m8_measure_level_statement,
    anc:m8_current_twenty_agent_split,
    anc:m8_current_twelve_agent_wave,
    anc:m8_current_eighteen_agent_wave,
    anc:m8_proof_wave_true_nodes}
  \begingroup\raggedright
  The remaining blockers should be reported with priority and field name.

  \emph{P0: real measure-localization blockers.}
  \begin{enumerate}
    \item \texttt{bulkIntegralLocalizes} should be filled through
      \texttt{BulkMeasureLocalizationTermFields.bulkIntegralLocalizes} or the
      \texttt{BulkIntegralLocalizationInput} constructor.  The real missing
      data are the chartwise bulk measure terms, the a.e. integrand
      replacement for the actual measure, integrability of every localized
      term, and the measure-box comparison saying that the measure-local term
      is the project-local box term.
    \item \texttt{boundaryMeasureIntegral} and
      \texttt{boundaryMeasureIntegral\_eq\_partitionSum} should be filled
      through \texttt{BoundaryMeasureLocalizationData} or the COV constructor.
      The real missing data are the ambient boundary measure and integrand, the
      a.e. indicator decomposition over selected boundary pieces, termwise
      integrability, \texttt{globalBoundaryIntegral = boundaryMeasureIntegral},
      and piece-integral equality with the selected partition terms.
  \end{enumerate}

  \emph{P1: geometric instances behind existing constructors.}
  \begin{enumerate}
    \item Artificial faces are fieldized by
      \texttt{ArtificialFaceResolvedData}.  The missing instance is either a
      real pairing of opposite selected artificial faces with equal unsigned
      terms and opposite signs, or a support-zero proof showing that every
      localized artificial face term vanishes.
    \item Target-image selection is fieldized on the boundary-chart side and
      has a global adapter.  The missing instance is the local-inverse/open-map
      construction producing target selected boxes, image equality, local
      bijection, and the alignment between target boxes, partition boxes, and
      selected boxes.
  \end{enumerate}

  \emph{P2: packaging and API alignment.}
  \begin{enumerate}
    \item Boundary reconstruction wrappers and COV wrappers already exist, but
      selected representatives still have to use the same term names as the
      assembly layer.
    \item Import hygiene remains part of the task: global adapters may import
      boundary-chart work, but \texttt{Stokes.HalfSpace} must not import
      \texttt{Stokes.Global}.
    \item The final M8 statement should remain a projection from existing M7
      data plus \texttt{M8MeasureLocalizationData}; do not add a new final
      constructor layer unless a field genuinely cannot be expressed there.
  \end{enumerate}
  \par\endgroup
\end{remark}

\begin{theorem}[M8 to manifold Stokes dependency ladder]
  \label{anc:m8_to_manifold_dependency_ladder}
  \uses{anc:m8_p0_p1_p2_blocker_ledger,
    anc:m8_measure_level_statement,
    thm:global_compact_support_stokes,
    thm:compact_manifold_stokes}
  \begingroup\raggedright
  The route from M8 to the usual manifold theorem should be split into the
  following assignable layers.
  \begin{enumerate}
    \item Local layer: half-space and boundary-chart local Stokes supply the
      project-local bulk and outward-first boundary terms.
    \item Geometry layer: compact support and chart sources select finite
      interior and boundary boxes, extended boxes, and target image boxes.
    \item Partition layer: a smooth finite active partition produces localized
      forms, support containment, smoothness on the chosen neighborhoods, and
      support-local exterior-derivative equality.
    \item Measure layer: the bulk and boundary measure-localization packages
      identify the genuine global integrals with the selected finite sums.
    \item Cancellation layer: artificial interior faces cancel, and real
      boundary terms are transported through oriented boundary chart changes.
    \item M8 assembly layer: \texttt{M8MeasureLocalizationData} plus the M7
      selected mixed input gives \texttt{m8GlobalStokes}.
    \item Compact-support statement: the selected partition construction fills
      the M8 input from a compactly supported chartwise smooth form.
    \item Compact manifold statement: compactness supplies the compact-support
      wrapper, yielding the ordinary compact oriented smooth
      manifold-with-boundary theorem.
    \item Mathlib comparison: if mathlib later provides a canonical manifold
      differential-form integral, prove that the project-local integral agrees
      with it.
  \end{enumerate}
  This ladder is the intended division of labor after M8: early workers should
  discharge orange geometric and measure instances; later workers should only
  assemble stable records into theorem statements.
  \par\endgroup
\end{theorem}

\endgroup
